\chapter{Captura de requisitos con musicólogas y formulación de preguntas para comprobar las capacidades de los modelos de lenguaje}
\label{cap:capitulo3}

\section{Captura de requisitos con musicólogas a través de entrevistas personales}
Para recopilar información sobre las necesidades, expectativas y posibles usos de la herramienta desarrollada, se realizaron entrevistas presenciales con musicólogas de la Facultad de Geografía e Historia de la Universidad Complutense de Madrid. Estas entrevistas permitieron obtener información cualitativa más detallada sobre los criterios de búsqueda, análisis y contextualización considerados relevantes en el ámbito de la musicología y el estudio de la música tradicional y folclórica.

Entre los principales intereses identificados destacó la necesidad de realizar búsquedas musicales avanzadas a partir de elementos concretos de las obras, como los íncipits musicales (los primeros compases instrumentales o vocales), estructuras melódicas, letras, instrumentación o características interpretativas. Asimismo, las investigadoras señalaron la importancia de disponer de información musicológica detallada, incluyendo aspectos como autoría, fecha, duración, tonalidad, tempo, rango vocal o presencia de grabaciones sonoras.

Otro de los aspectos recurrentes fue el interés por analizar la evolución y transmisión de las obras a lo largo del tiempo, especialmente en el contexto de la música tradicional y folclórica. En este sentido, se destacó la relevancia de estudiar las variaciones entre versiones, los procesos de creación y transmisión oral, así como las relaciones existentes entre músicas de distintas regiones, comunidades y contextos culturales.

Las entrevistas también pusieron de manifiesto la importancia de incorporar funcionalidades orientadas a la comparación y análisis relacional de obras musicales, tales como la identificación de patrones compartidos, la frecuencia de aparición de determinados íncipits o leitmotivs musicales dentro de un corpus o la comparación simultánea de canciones dentro de una misma interfaz.

Asimismo, las participantes señalaron la necesidad de integrar información procedente de diferentes catálogos y bases de datos especializadas (como BNE, ISNI, Wikidata o SGAE), incluyendo variantes de nombres y pseudónimos de autores e intérpretes con el fin de facilitar procesos de recuperación y contextualización de la información musical.

Finalmente, las entrevistas evidenciaron el valor de conservar metadatos relacionados con la recopilación y procedencia de las obras, tales como el lugar de recuperación, el contexto de documentación o los datos asociados al proceso de recogida de la información, aspectos especialmente relevantes en investigaciones vinculadas al patrimonio musical y la tradición oral.

\section{Captura de requisitos con musicólogas a través de formularios}
Además de las entrevistas realizadas de manera presencial, se diseñó un formulario estructurado dirigido a potenciales usuarias pertenecientes al ámbito académico e investigador. El objetivo principal de este instrumento fue identificar los perfiles de los participantes, conocer los tipos de consultas que realizarían con la herramienta y comprender el uso previsto de los resultados obtenidos.
El cuestionario fue creado y difundido mediante Google Forms, lo que permitió una distribución sencilla y accesible, así como la recopilación organizada de las respuestas.

El cuestionario incluyó, en primer lugar, una pregunta destinada a caracterizar el perfil académico de los participantes, ofreciendo las categorías de investigador/a, estudiante, profesor/a u otras situaciones profesionales relacionadas. Asimismo, se solicitó información sobre el área de investigación de cada encuestada con el fin de contextualizar sus respuestas y analizar posibles diferencias entre disciplinas.

A continuación se incorporaron dos preguntas abiertas orientadas a identificar las consultas que las usuarias podrían considerar relevantes a realizar mediante la herramienta y a documentar el uso que las académicas darían a los datos proporcionados por el sistema tras realizar dicha consulta. 
La información recopilada permitió, por una parte, recoger ejemplos concretos de necesidades de búsqueda, análisis o recuperación de información, así como otros posibles escenarios de uso no previstos inicialmente por los investigadores y, por otra parte, explorar el potencial impacto de la herramienta en actividades de investigación, docencia, análisis, documentación o generación de nuevo conocimiento.

Las respuestas obtenidas mediante el formulario permitieron identificar distintos intereses y necesidades relacionadas con la investigación musicológica y el desarrollo de herramientas de recuperación y análisis de información musical. Las participantes pertenecían principalmente al ámbito investigador y docente, con especialización en áreas como la música latinoamericana de los siglos XIX y XX, la producción musical y el teatro lírico, así como la musicología y los estudios sobre flamenco y folclore.

Uno de los aspectos más destacados fue el interés por disponer de funcionalidades de búsqueda avanzadas que permitan localizar distintas versiones de una misma obra, clasificar repertorios por estilos o estructuras musicales y acceder a diferentes interpretaciones y grabaciones relacionadas con un mismo material musical. Asimismo, varias respuestas señalaron la importancia de incorporar conexiones con otros recursos externos, como archivos, bibliotecas o bases de datos especializadas, con el fin de contextualizar las obras y establecer relaciones entre distintos repertorios musicales y literarios.

En el ámbito de la música tradicional y el folclore, las respuestas pusieron de manifiesto la relevancia de incluir criterios de búsqueda relacionados con aspectos temáticos, geográficos y socioculturales. Entre ellos destacan la posibilidad de identificar referencias a lugares concretos, analizar temáticas recurrentes en las letras o estudiar variaciones lingüísticas y culturales en función de la región, el idioma o la época histórica.

Asimismo, las musicólogas señalaron el interés de utilizar los datos recuperados para la construcción de bases de datos propias, el desarrollo de proyectos de investigación, la elaboración de materiales docentes y la creación de repertorios especializados. También se destacó el potencial de la herramienta para generar visualizaciones y mapas interactivos que permitan representar la distribución geográfica y la evolución de determinadas tradiciones musicales.

Otro aspecto recurrente fue la importancia atribuida a la documentación del proceso de recopilación y transmisión de la música tradicional. Se consideró relevante incorporar información sobre la labor desarrollada por folcloristas y etnomusicólogos, incluyendo datos sobre las regiones estudiadas, los repertorios recopilados y las metodologías de recogida empleadas.

En conjunto, los resultados del formulario evidencian la necesidad de una herramienta flexible y adaptada, capaz de integrar información procedente de múltiples fuentes y de ofrecer funcionalidades avanzadas de búsqueda, análisis y contextualización musical. Las respuestas obtenidas refuerzan además la importancia de abordar la música no únicamente como un objeto textual o sonoro aislado, sino también como un fenómeno cultural, histórico y social.

\section{Formulación de preguntas para comprobar las capacidades de los modelos de lenguaje}
Como parte del diseño experimental inicial, se elaboró un conjunto de preguntas teórico-musicales orientadas a evaluar las capacidades potenciales de los modelos de lenguaje pensados para la herramienta en distintos ámbitos de la musicología, el análisis musical y la recuperación avanzada de información musical. El objetivo de este proceso fue definir un repertorio amplio y representativo de consultas que permitiera cubrir tanto aspectos estrictamente musicales —como ritmo, tonalidad, métrica o melodía— como dimensiones documentales, textuales, históricas y técnicas relacionadas con los formatos de codificación musical utilizados en el corpus.

Las preguntas fueron formuladas con distintos niveles de complejidad y especificidad, combinando criterios musicales, metadatos descriptivos, análisis lírico y aspectos estructurales derivados de formatos como MEI, MusicXML o MSCZ. De este modo, se buscó evaluar la viabilidad de consultas complejas capaces de relacionar múltiples dimensiones analíticas dentro de un mismo entorno de recuperación de información.

Para facilitar su organización y análisis, las preguntas fueron agrupadas en diferentes bloques temáticos en función del ámbito musicológico o técnico abordado.

\begin{itemize}
\item \textbf{Ritmo, métrica y tempo.} Preguntas centradas en la acentuación, figuras rítmicas, síncopas, polirritmias, velocidad (BPM) y la relación matemática/temporal de las notas dentro del compás. Preguntas incluidas:
    \begin{itemize}
    \item \textit{Identifica piezas en 2/4 que contengan síncopas ligadas entre el segundo pulso de un compás y el primero del siguiente de forma sistemática}
    \item \textit{Compara la densidad de notas por compás entre una pieza de 4/4 y una de 2/4 del dataset; ¿cuál tiene mayor promedio de eventos musicales?}
    \item \textit{Identifica si existe algún uso de polirritmia en los archivos MEI}
    \item \textit{Identifica si existe algún patrón de silencios (<rest />) que se repita sistemáticamente cada cuatro compases en las piezas de métrica ternaria}
    \end{itemize}
\item \textbf{Análisis melódico y armónico.} Consultas orientadas a los intervalos musicales, tesituras, movimientos melódicos (ascendentes/descendentes), alteraciones accidentales (sostenidos, bemoles y becuadros), tonalidades y grados de resolución (tónica/dominante). Preguntas incluidas:
    \begin{itemize}
    \item \textit{Localiza todas las notas que tengan bemoles accidentales y dime en qué compases aparecen en los archivos MEI encontrados}
    \item \textit{¿Cuál es el intervalo melódico más frecuente entre la última nota de un compás y la primera del siguiente en las piezas de modo menor?}
    \item \textit{Identifica piezas donde la nota final sea la dominante (quinto grado) en lugar de la tónica, sugiriendo una estructura abierta}
    \end{itemize}
\item \textbf{Relación métrica, lírica y texto.} Preguntas híbridas que analizan cómo interactúa el texto (la lírica, el número de sílabas, palabras clave) con los aspectos musicales del tempo, la altura de las notas o la duración rítmica. Preguntas incluidas:
    \begin{itemize}
    \item \textit{Identifica canciones donde el autor sea una mujer, pero el contenido lírico esté escrito en voz masculina}
    \item \textit{Filtra piezas que usen la palabra "estrella" en la lírica y determina si su melodía es predominantemente ascendente o descendente}
    \item \textit{Compara el uso de la sílaba "duer-me" en los distintos archivos del dataset: ¿se asigna siempre a la misma duración de nota?}
    \item \textit{Determina si las canciones con tempo superior a 120 BPM tienden a tener letras con palabras de menor longitud silábica}
    \item \textit{Busca el término "cie-lo" en las líricas y analiza si la nota asignada tiene un oct (octava) superior a 4}
    \end{itemize}
\item \textbf{Contexto sociocultural, temático y geográfico.} Preguntas complejas que cruzan la etnomusicología o el análisis musical comparativo con datos del entorno (temas como "Nanas" o "Trabajo", regiones geográficas o fechas de recolección históricas). Preguntas incluidas:
    \begin{itemize}
    \item \textit{Identifica piezas del dataset compuestas entre 1750 y 1790 que utilicen una hemiolia sistemática en la transición del estribillo al verso, pero solo en aquellas cuya tonalidad sea menor armónica y que no superen los 90 BPM. Explica cómo la letra refuerza ese cambio de acentuación rítmica}
    \item \textit{Busca ejemplos de anacrusa en piezas que estén etiquetadas como "Nanas"}
    \item \textit{Lista canciones que compartan la misma keySig de un bemol (Fa mayor/Re menor) pero que procedan de regiones geográficas lejanas entre sí}
    \item \textit{Encuentra piezas recolectadas antes de 1950 que traten el tema de "la muerte" y que utilicen exclusivamente un rango melódico de sexta menor}
    \item \textit{¿Qué diferencias de tesitura (nota más aguda vs. más grave) existen entre las versiones de "Nanas" de los distintos países presentes en el dataset?}
    \item \textit{Busca canciones de "trabajo" (agrícolas) que compartan la misma estructura de intervalos que las "Nanas" del dataset}
    \item \textit{Encuentra canciones que utilicen leitmotivs rítmicos idénticos en la región andina y en la península ibérica}
    \end{itemize}
\item \textbf{Estructura analítica de datos y formatos técnicos.} Este bloque engloba las preguntas que se centran puramente en el código o metadatos de los archivos musicales (MEI, MusicXML, XML, MSCZ), discrepancias de sintaxis musical y atributos técnicos de renderizado (como plicas, volúmenes MIDI o transcripciones). Preguntas incluidas:
    \begin{itemize}
    \item \textit{Encuentra canciones donde el valor de dur.ppq cambie a mitad de la sección, indicando un cambio de resolución rítmica}
    \item \textit{Localiza obras que tengan compases donde la suma de las duraciones de las etiquetas <note> no coincida con el meterSig declarado en el staffDef}
    \item \textit{¿Qué piezas utilizan valores irregulares (por ejemplo, tresillos) no declarados explícitamente en un <time-modification> del MusicXML?}
    \item \textit{Identifica el uso de la etiqueta <sb /> (system break) y determina si coincide siempre con el final de una frase musical lógica}
    \item \textit{¿Qué archivos contienen metadatos en la etiqueta <appInfo> que indiquen que fueron transcritos desde MusicXML a MEI usando Verovio?}
    \item \textit{Busca discrepancias entre el beat-unit del metrónomo y el beat-type de la armadura de tiempo en los archivos .xml}
    \item \textit{Encuentra notas que tengan el atributo stem.dir="down" a pesar de estar por debajo de la tercera línea del pentagrama}
    \item \textit{Identifica piezas que utilicen la etiqueta <instrDef> con un midi.volume inferior al 80\%}
    \item \textit{¿Qué archivos MSCZ presentan metadatos de "Quality Control" firmados por una persona específica?}
    \end{itemize}
\end{itemize}